\begin{abstract}


\draftcomment[Karen]{ What follows is our current state of play from my perspective, what can we add in time?  What can you add from your own work?  What we want in the end is to draw from all our work and present a common perspective.}

This paper reports a reproduction and deployment-oriented adaptation of the MuJoCo Playground OP3 joystick locomotion task for the ROBOTIS OP3 humanoid. The work is not presented as a novel reinforcement-learning algorithm. Instead, it documents the engineering path from a reference MuJoCo Playground policy to a STRIDE training and deployment pipeline: actuator-order alignment with ROBOTIS motor IDs, Brax PPO training in MJX, ONNX export, ROS~2 deployment, MuJoCo-to-Webots sim-to-sim testing through the ROBOTIS OP3 ROS~2 framework, and measurement-driven sim-to-real diagnosis. The main contribution is a set of practical lessons about why a policy that walks in simulation may fail on hardware: observation and reward conventions must match the reference implementation exactly, the deployment control path can accidentally filter policy commands, and the default OP3 actuator stiffness in the simulator is substantially softer than measured hardware behaviour. We argue that this kind of faithful reproduction and model-fidelity audit is a necessary precursor to credible sim-to-real claims for small humanoid locomotion.
\end{abstract}

\noindent\textbf{Keywords:} ROBOTIS OP3; MuJoCo Playground; Webots; MJX; Brax PPO; humanoid locomotion; sim-to-sim; sim-to-real; reproducibility.

\section{Introduction}
\label{sec:introduction}

Reinforcement learning can now produce visually convincing locomotion policies for small humanoid robots in high-throughput physics simulators. However, the practical value of these policies depends on whether the simulated task, the exported policy, and the real robot control path describe the same closed-loop system.

This paper presents a reproduction and adaptation of the MuJoCo Playground OP3 joystick locomotion environment \citep{Todorov2012-mujoco-physics-engine-model-based-control,Freeman2021-brax-differentiable-physics-engine-large-scale-rigid-body-simulation,Zakka2025-demonstrating-mujoco-playground,GoogleDeepMind2025-mujoco-playground}. The aim is deliberately modest: reproduce the reference OP3 walking result, adapt it to a ROBOTIS-motor-ordered STRIDE training pipeline, and identify the main sim-to-real gaps before making claims about robust physical deployment.

\textbf{Draft positioning.} This is not new algorithmic work. The technical contribution is an engineering and measurement report: what had to match, what initially did not match, and which simulator assumptions were contradicted by measurements from the physical OP3.

\section{Background and Related Work}
\label{sec:background}

\subsection{MuJoCo, MJX, and Playground}

MuJoCo provides the rigid-body simulation substrate and MJX provides the JAX-compatible accelerated execution path used for high-throughput policy training. MuJoCo Playground builds on this stack as a canonical open-source suite of robot-learning environments for locomotion, manipulation, and vision-based control \citep{Todorov2012-mujoco-physics-engine-model-based-control,Zakka2025-demonstrating-mujoco-playground,GoogleDeepMind2025-mujoco-playground-gpu-accelerated-robot-learning-sim-to-real-transfer}. The raison d'etre of Playground is to make sim-to-real robot learning faster to start, faster to iterate, and easier to reproduce. The RSS paper presents it as a fully open-source MJX-based framework whose integrated stack combines GPU-accelerated physics, on-device batch rendering, ready-made training environments, notebooks, hyperparameters, training curves, and deployment examples. The intended effect is to move the user's effort away from repeatedly rebuilding simulator and training infrastructure, and toward encoding desired robot behaviour, training policies on a single GPU, and testing whether those policies transfer to hardware \citep{Zakka2025-demonstrating-mujoco-playground}.

For this study, the important point is that ROBOTIS OP3 is included as a first-class humanoid joystick environment alongside Berkeley Humanoid, Unitree H1/G1, and Booster T1, sharing the same commanded velocity abstraction $(v_x, v_y, \omega_z)$. The paper reports rapid zero-shot deployment across six real robot platforms, and the RSS~2025 demo page records live demonstrations on Unitree Go1, Unitree G1, Booster T1, and LEAP hand \citep{Zakka2025-demonstrating-mujoco-playground,MuJoCoPlaygroundContributors2025-mujoco-playground-rss-2025}. OP3 is therefore supported by the canonical library, but not reported as a demonstrated zero-shot hardware platform; this gap is the focus of our deployment audit.

The reference implementation defines the observation layout, reward terms, command sampling, PD gain setup, history stacking, and PPO training recipe used as the baseline.

\subsection{Brax PPO for locomotion}

Explain why the working pipeline uses Brax PPO rather than a custom PPO implementation: scan-based rollout collection, vectorised environments, observation normalisation, and much higher sample throughput.

\subsection{ROBOTIS OP3 deployment context}

Briefly describe the OP3: 20 XM430-W350 Dynamixel joints, position-control deployment, 50~Hz RL policy inference, ROS~2 control path, and the practical constraints imposed by the ROBOTIS framework.

\section{System Overview}
\label{sec:system}

\subsection{Reference and STRIDE models}

The reference model is the MuJoCo Playground OP3 joystick locomotion task \citep{Zakka2025-demonstrating-mujoco-playground,GoogleDeepMind2025-mujoco-playground-gpu-accelerated-robot-learning-sim-to-real-transfer}. The STRIDE model keeps the same broad training formulation but reorders actuators to match ROBOTIS motor IDs 1--20, so the policy output vector maps directly to real robot motor commands.

\subsection{Training pipeline}

The current pipeline is:
\[
\text{MuJoCo MJCF} \rightarrow \text{MJX} \rightarrow \text{Brax PPO} \rightarrow \text{Orbax checkpoint} \rightarrow \text{ONNX} \rightarrow \text{ROS 2 deployment}.
\]

Important training parameters to report:
\begin{itemize}
  \item 50~Hz control, 0.004~s simulation timestep, 5 simulation steps per action.
  \item 2048 parallel environments in the usual training configuration.
  \item 147-dimensional observation: 49-dimensional frame stacked over three history frames.
  \item 20-dimensional action, squashed to $[-1,1]$ and scaled by 0.3~rad around the default pose.
  \item Policy network: four hidden layers of 128 units; critic: five hidden layers of 256 units.
\end{itemize}

\subsection{Observation and command conventions}

The policy observes body-frame angular velocity, upvector, commanded velocity, joint positions relative to the default pose, previous action, and two previous observation frames. The commanded joystick velocity is a continuous triple $(v_x, v_y, \omega_z)$ sampled during training and supplied by the deployment command interface at runtime.

\section{Reproducing Playground Walking}
\label{sec:reproduction}

\subsection{Custom PPO failure mode}

Describe the initial custom PPO attempts: stable standing, poor walking, low throughput, missing details relative to Playground. Use this as a reproducibility lesson rather than as a negative result against PPO.

\subsection{Reference-matching fixes}

The important fixes were:
\begin{itemize}
  \item reward scaled by control timestep and clipped to be non-negative;
  \item body-frame rather than world-frame velocity tracking;
  \item observation normalisation from the start of training;
  \item three-frame observation history;
  \item Playground-compatible command ranges and resampling;
  \item PD gain setup through MuJoCo actuator parameters;
  \item Brax scan-based rollout and update loop.
\end{itemize}

\subsection{Training result}

Insert training curves and short rollout snapshots here. Suggested figure: reward versus environment steps comparing custom PPO, Playground reference, and STRIDE/Brax PPO.

\section{Policy Export and Deployment}
\label{sec:deployment}

\subsection{ONNX export}

The exported policy bakes observation normalisation into the graph and uses the deterministic action $a=\tanh(\mu)$, discarding the stochastic log-standard-deviation outputs used during training.

\subsection{ROS 2 control path}

Document the 50~Hz policy node, observation construction from IMU quaternion, joint states, command velocity, last action, and history stacking.

\subsection{Deployment control filtering}

The ROBOTIS \texttt{direct\_control\_module} can strongly low-pass filter position commands through minimum-jerk smoothing. For an RL policy emitting commands every 20~ms, moving-time and moving-angle parameters determine whether the hardware receives near-passthrough commands or heavily filtered trajectories.

\section{Sim-to-Sim Transfer: MuJoCo to Webots}
\label{sec:sim2sim}

Before treating hardware behaviour as a sim-to-real result, we use Webots as an intermediate simulator because it exercises the policy through the ROBOTIS OP3 ROS~2 framework rather than through the MuJoCo training loop. This tests whether the exported policy, observation builder, motor ordering, command timing, and ROS~2 control interfaces remain consistent outside the original MJX environment.

\subsection{Purpose of the Webots stage}

The Webots stage is not intended to replace MuJoCo as the training environment. Its role is a deployment rehearsal: the same ONNX policy and ROS~2 policy node used for the physical OP3 are connected to a simulated OP3 in Webots. Disagreements between MuJoCo and Webots are therefore useful even when Webots is not assumed to be more physically accurate. They expose assumptions that are hidden when policy inference and physics stepping occur inside the same simulator.

\subsection{ROBOTIS OP3 ROS 2 integration}

Document the ROS~2 graph used for Webots testing: policy inference node, joint-state source, IMU source, command velocity input, and position command output. The important check is that the Webots bridge presents the same interface expected by the hardware-facing ROBOTIS OP3 modules, including joint names, motor order, units, update rate, and command semantics.

\subsection{Quantities to compare}

Report MuJoCo and Webots rollouts under matched joystick commands. Useful comparisons include commanded versus achieved base velocity, torso pitch and roll, foot contact timing, joint position tracking, action saturation, and failure modes such as foot scuffing, collapse, or persistent turning bias. This section should make clear which differences are simulator-dependent and which differences indicate an error in the exported policy or ROS~2 deployment path.

\subsection{Expected outcomes}

The expected contribution of this stage is diagnostic rather than a claim that Webots validates the policy. A successful sim-to-sim test shows that the policy can be driven through the deployment stack in a second simulator without relying on MuJoCo-specific implementation details. A failed test narrows the search space before hardware experiments, especially for issues involving motor ordering, coordinate frames, observation history, policy timing, and command filtering.

\section{Sim-to-Real Fidelity Findings}
\label{sec:sim2real}

\subsection{Motor stiffness mismatch}

The MuJoCo Playground default OP3 stiffness, $K_p=21.1$, is too soft for the measured physical robot. In static walk-ready recordings, the real OP3 held all joints within roughly 0.1--0.2 degrees over 5 seconds, while the default simulator showed substantially larger yielding. Increasing simulated stiffness towards $K_p=100$ reduced the discrepancy but did not fully match the real robot.

\subsection{Trajectory-shape mismatch}

Action-page replay initially showed confusing dynamic errors: a softer simulated controller sometimes tracked better than a stiffer, more realistic controller. The reason was an input mismatch. The simulator was replaying linear interpolation between keyframes, while the ROBOTIS action module uses a trapezoidal velocity profile. Porting the trapezoidal generator greatly reduced dynamic tracking error.

\subsection{Foot friction and spikes}

The real OP3 feet include protruding screws that behave like spikes on grass. A policy trained with moderate Coulomb friction may learn to slide or pivot in ways that fail on hardware. Future training should randomise friction, penalise foot slip more strongly, and represent screw-tip contacts or equivalent high-friction point contacts.

\section{Discussion}
\label{sec:discussion}

\subsection{Reproduction as scientific work}

The core lesson is that reproducing a locomotion result requires reproducing the full system: simulator settings, observation conventions, reward scaling, control rates, export semantics, and hardware command filtering. Small mismatches can dominate the observed behaviour.

\subsection{What is original here}

The original contribution is not a new PPO variant or locomotion reward. It is the OP3-specific audit trail: a deployment-aligned actuator order, measured simulator/hardware discrepancies, diagnosis of command-filtering effects, and a concrete roadmap for retraining against a more faithful model.

\subsection{Limitations}

Current limitations to state honestly:
\begin{itemize}
  \item the current policy was trained against a motor model now known to be too soft;
  \item robust walking on the physical robot is not yet the claimed endpoint;
  \item static stiffness measurements do not fully identify dynamic actuator behaviour;
  \item battery voltage, torso mass, foot-ground interaction, and control latency remain partly unmodelled.
\end{itemize}

\section{Conclusion}
\label{sec:conclusion}

Conclude that the work turns a MuJoCo Playground reference task into a reproducible OP3 training and deployment pipeline, and that the main scientific value lies in exposing the hidden assumptions that must be made explicit before sim-to-real locomotion claims are credible.

\section*{Acknowledgements}

Placeholder for STRIDE funding from SPECS [ ? ] and acknowledgements [conversations with Abolfazi ? ].
